\section{Introduction}
Two-dimensional platform games provide a useful environment for studying game development because they combine real-time input, animation, physics, collision detection, artificial intelligence, resource management, and user-interface design. Although their basic objective is easy to understand, implementing a reliable platform game requires many interconnected systems to operate consistently. Character movement must respond smoothly to player input, collisions must be resolved without allowing entities to pass through the environment, enemies must react appropriately, and visual animations must remain synchronized with the underlying gameplay state.

The Mario Game Project was developed as a custom C++ platform game inspired by the mechanics of the Super Mario series. The game allows players to control Mario, Luigi, or Flash. Each hero shares a common movement and interaction framework while providing different textures, animations, and projectile-based special abilities. Heroes can collect mushrooms to enter a larger form and flowers to unlock an advanced form, such as Fire Mario or ThunderFlash. Other gameplay elements include collectible coins, temporary invincibility, interactive blocks, pipes, environmental obstacles, and multiple enemy types.

A major objective of the project is to apply object-oriented programming principles to a practical game system. Common entity behavior is represented through inheritance and polymorphism, while specialized classes implement the unique behavior of heroes, enemies, items, blocks, and projectiles. Character forms and movement conditions are modeled using separate form and state classes. This approach allows transitions such as growing, shrinking, jumping, crouching, taking damage, and celebrating victory to be managed without placing all behavior inside a single character class. Factory classes are also used to create heroes, enemies, blocks, and items from level data.

The project separates the gameplay world from presentation and application-level state management. Dedicated systems handle physics, entity interactions, level construction, runtime updates, animation, sound, HUD information, menus, pause behavior, victory, defeat, and transitions between screens. Levels are loaded from Tiled map files, enabling terrain, triggers, enemy spawners, interactive objects, pipe routes, and goal regions to be configured outside the source code. A JSON-based save mechanism records important information such as the active level, hero status, enemy data, modified blocks, score, coins, lives, and remaining time.

This report presents the design and implementation of the Mario Game Project, with particular attention to its software architecture, gameplay systems, object-oriented design, enemy behavior, physics, animation, user interface, and persistence mechanism. It also discusses integration challenges encountered during development and identifies areas for future improvement.